\providecommand{\NTests}{334}\providecommand{\NRuns}{3}\providecommand{\NClean}{3}

\providecommand{\NFR}{68}\providecommand{\NAuto}{52}\providecommand{\NAutoUI}{10}\providecommand{\NPartial}{2}\providecommand{\NInspect}{2}\providecommand{\NManual}{2}

\providecommand{\CollN}{50}\providecommand{\CollDel}{92.6}\providecommand{\CollDelMin}{76}\providecommand{\CollDelMax}{100}\providecommand{\CollRecMean}{8.1}\providecommand{\CollRecMax}{10.25}\providecommand{\CollRecN}{50}\providecommand{\CollRecOK}{23}\providecommand{\CollConvMean}{9.6}\providecommand{\CollConvMax}{15.2}

\providecommand{\HopW}{99.2}\providecommand{\HopH}{57.4}\providecommand{\HopWins}{20}\providecommand{\HopRuns}{20}

\providecommand{\SosAdv}{4.4}

\providecommand{\NDevChecks}{21}\providecommand{\NDevPass}{21}\providecommand{\NWebChecks}{16}\providecommand{\NWebPass}{16}

\section*{Abstract}
\addcontentsline{toc}{section}{Abstract}

\textbf{Problem.} An earthquake, fire or collapse can take down the towers and internet links that every popular messaging application needs, although the phones of the people trapped inside can often still reach one another. Whatever connectivity survives helps only the person who holds it.

\textbf{Motivation and objectives.} Rescue is time-critical, so the aim is to let one person's connectivity serve everyone nearby and to let an SOS overtake routine traffic at every hop: a self-organising mesh, routing weighted by measured link quality, reliable delivery with a hold-and-forward fallback, encrypted and signed messages, and measured performance.

\textbf{Proposed solution.} In EcoMesh every device both sends and relays. A message moves device by device until it reaches a gateway with an outside link or a rescue device, or is held until a path appears. No internet, tower or server is needed.

\textbf{Technologies and engineering concepts.} Python \texttt{asyncio} (one event-loop thread per node); a 52-byte binary header of our own design; link-state gossip; Dijkstra over the cost $1+5(1-q)+3f$; a priority queue at every hop; ACK with retry and re-routing; AES-256-GCM and Ed25519; a SQLite log per node; FastAPI and WebSocket; React; and a transport interface that runs one engine on a simulator, on loopback and across laptops.

\textbf{Key implementation features.} Discovery with measured link quality; link-state flooding controlled by sequence numbers; duplicate suppression; a bounded hold queue that never drops an SOS for routine traffic; signatures that survive relaying; and one device app that becomes a survivor screen (hold-to-send SOS with nine request types, location, voice input, ``I am safe'', the rescue reply), a gateway view or a rescue console (requests with their route, replies, filters, history, CSV export, live map and activity panel).

\textbf{Testing methodology.} \NTests\ automated tests (unit, integration, boundary, fuzz, replay, concurrency, a three-node mesh over TCP), a test mapping for all 68 requirements, seven seeded simulator scenarios, loss and scale sweeps to 65 nodes, five live processes over real sockets, and scripted browser checks of the device screens (\NDevPass\ of \NDevChecks\ pass).

\textbf{Results.} A 40-node mesh at 20\,\% loss with injected failures delivered \CollDel\,\% of messages (50 runs). On identical seeds weighted routing delivered \HopW\,\% against \HopH\,\% for hop count. Under load an SOS had a \SosAdv$\times$ lower median latency than routine traffic. On five live processes an SOS arrived in a median 32\,ms, a sender re-routed around a failed relay in a median 4.1\,s, and every message held during a gateway outage was delivered. Convergence and recovery met their targets on average but not in every run; at 20\,\% loss delivery stayed above 90\,\% to 50 nodes and fell to 83.1\,\% at 65.

\textbf{Practical applicability and limitations.} EcoMesh can serve today as an institutional safety network on laptops. The radio is emulated by a LAN restricted to configured neighbours, so it is not yet phone to phone; a Bluetooth or Wi-Fi Direct transport, adaptive timers, a usability study and per-pair keys are the main future work.
